\DocumentMetadata{
  pdfversion=2.0,
  pdfstandard=ua-2,
  lang=en-US,
  tagging=on,
  tagging-setup={math/setup=mathml-SE,tabsorder=structure}
}
\documentclass[11pt,letterpaper]{article}
\usepackage[margin=0.68in,includefoot]{geometry}
\usepackage{newcomputermodern}
\usepackage{amsmath,amscd,mathtools}
\usepackage{tikz,adjustbox}
\providecommand{\square}{\mdlgwhtsquare}
\usepackage[hidelinks]{hyperref}
\hypersetup{
  pdftitle={Analysis First-Year Exam, January 2023, Part 2},
  pdfauthor={University of Florida Department of Mathematics},
  pdfsubject={Graduate mathematics examination},
  pdfdisplaydoctitle=true,
  bookmarksopen=true
}
\setcounter{secnumdepth}{0}
\setlength{\parindent}{0pt}
\setlength{\parskip}{0.28em}
\setlength{\emergencystretch}{3em}
\setlength{\leftmargini}{2.15em}
\setlength{\leftmarginii}{2.65em}
\setlength{\leftmarginiii}{3em}
\renewcommand{\labelenumi}{\textbf{\theenumi.}}
\pagestyle{empty}
\raggedbottom
\allowdisplaybreaks
\newenvironment{examproblems}[1]{%
  \begin{enumerate}%
  \setcounter{enumi}{#1}%
  \setlength{\topsep}{0.35em}%
  \setlength{\partopsep}{0pt}%
  \setlength{\itemsep}{0.48em}%
  \setlength{\parsep}{0.18em}%
}{\end{enumerate}}
\newenvironment{examsubparts}{%
  \begin{enumerate}%
  \setlength{\topsep}{0.18em}%
  \setlength{\partopsep}{0pt}%
  \setlength{\itemsep}{0.16em}%
  \setlength{\parsep}{0.1em}%
}{\end{enumerate}}
\newenvironment{examinstructions}{%
  \par\begingroup\itshape
}{%
  \par\endgroup\vspace{0.18em}
}
\makeatletter
\renewcommand\section{\@startsection{section}{1}{0pt}%
  {-1.25ex plus -.25ex minus -.1ex}%
  {0.55ex plus .1ex}%
  {\normalfont\large\bfseries}}
\renewcommand{\@maketitle}{%
  \newpage\null\vskip -1em%
  {\footnotesize\bfseries
    \href{https://www.ufl.edu/}{University of Florida}, \href{https://math.ufl.edu/}{Department of Mathematics}\par}%
  \vskip -0.5em%
  \tagstructbegin{tag=Title}%
    {\LARGE\bfseries\href{https://gma.math.ufl.edu/exams/analysis-first-year/}{Analysis First-Year Exam}, January 2023, Part 2\par}%
  \tagstructend%
  \vskip 0.25em%
  \tagmcbegin{artifact}%
    \hrule height 1pt%
  \tagmcend%
  \vskip 1em%
}
\makeatother
\title{Analysis First-Year Exam, January 2023, Part 2}
\author{}
\date{}
\begin{document}
\maketitle
\thispagestyle{empty}
\begin{examinstructions}
Answer FOUR questions in detail. State carefully any results used without proof.
\end{examinstructions}
\begin{examproblems}{0}
\item
Let~\(f\) be a real-valued function on \([0,1]\) and consider the statements below. Prove that (ii) implies (i); and show by example that (i)~\(\Rightarrow\) (ii) can fail.
\begin{examsubparts}
\item[\textnormal{(i)}]
for each positive integer~\(n\) the function~\(f\) is Riemann-integrable on \([1/n,1]\) and the sequence of integrals \(\int_{1/n}^1 f\) converges to a real number as \(n\to\infty\);
\item[\textnormal{(ii)}]
the function~\(f\) is Riemann-integrable on \([0,1]\).
\end{examsubparts}
\item
Let~\(X\) be a metric space on which \((f_n)\) and \((g_n)\) are uniformly convergent sequences of continuous real-valued functions.
\begin{examsubparts}
\item[\textnormal{(i)}]
Prove that if~\(X\) is compact then \((f_ng_n)\) is uniformly convergent.
\item[\textnormal{(ii)}]
Show that the conclusion in (i) can fail when~\(X\) is not compact.
\end{examsubparts}
\item
Let \(f:[0,1]\times[0,1]\to\mathbb R\) be continuous. By first considering polynomials in two variables, prove that 
\[
\int_0^1\left(\int_0^1 f(x,y)\,dx\right)dy=\int_0^1\left(\int_0^1 f(x,y)\,dy\right)dx.
\]
\item
Prove that if \((f_n)\) is a sequence of measurable functions such that \(f_n\to f\) on some measure space, then~\(f\) is measurable. Hence, or otherwise, show that if \(f:\mathbb R\to\mathbb R\) is differentiable then its derivative \(f'\) is Borel measurable.
\item
In each of the following cases, give (if possible and with justification) a sequence \((f_n)\) of Lebesgue integrable functions from \(\mathbb R\) to \(\mathbb R\) such that, where~\(\lambda\) denotes Lebesgue measure:
\begin{examsubparts}
\item[\textnormal{(i)}]
\(f_n\to0\) pointwise but \(\int f_n\,d\lambda\to1\) as \(n\to\infty\);
\item[\textnormal{(ii)}]
\(f_n\to0\) pointwise but \(\int f_n\,d\lambda\to\infty\) as \(n\to\infty\).
\end{examsubparts}
\end{examproblems}
\end{document}
